% Generated draft: requires mathematical and source review.
\documentclass[11pt]{article}
\usepackage[margin=25mm]{geometry}
\usepackage{fontspec}
\setmainfont{Latin Modern Roman}
\usepackage{amsmath,amssymb,mathtools,unicode-math}
\setmathfont{Latin Modern Math}
\usepackage{enumitem,needspace,xurl,hyperref}
\setlength{\parindent}{0pt}
\setlength{\parskip}{5pt}
\setlength{\emergencystretch}{3em}
\newcommand{\N}{\mathbb{N}}
\newcommand{\Z}{\mathbb{Z}}
\newcommand{\Q}{\mathbb{Q}}
\newcommand{\R}{\mathbb{R}}
\newcommand{\C}{\mathbb{C}}
\newcommand{\F}{\mathbb{F}}
\newcommand{\A}{\mathbb{A}}
\newcommand{\PP}{\mathbb{P}}
\newcommand{\E}{\mathbb{E}}
\newcommand{\eps}{\varepsilon}
\newcommand{\dd}{\,\mathrm{d}}
\begin{document}
\section*{UnsolvedMath ID: 20001117}
\textbf{Source problem code:} catalogue entry
\subsection{Definitions and mathematical statement}
Let $A$ be a finite nonempty subset of an arbitrary group $G$, with no commutativity, symmetry, or identity-inclusion assumption. Write $A^j=\{a_1\cdots a_j:a_i\in A\}$. T. Tao's original question asks whether small doubling forces small tripling, quantitatively whether a bound $|A^2|\leq K|A|$ controls $|A^3|/|A|$ by a function of $K$ alone. The source records that Ben Green and Mei-Chu Chang pointed out this is false, using $A=H\cup\{x\}$ with $H$ a finite subgroup and $x$ outside its normalizer. Here $|A^2|\leq3|H|+1$, whereas $A^3$ contains $HxH$, whose cardinality is $|H|[H:H\cap xHx^{-1}]$. A quantitative unbounded-growth family needs large conjugate index; non-normality alone is not a numerical growth estimate.

The source instead asks for a large subset with small tripling, and more ambitiously for inverse structure, especially under $|A^2|<2|A|$. The supplied aid reports the explicit corrected subset theorem:
\[
 \forall K\geq1:\quad |A^2|\leq K|A|
 \ \Longrightarrow\ \exists A'\subseteq A:\quad
 |A'|\geq |A|/K,\qquad |(A')^3|\leq K^3|A'|.
\]
The conclusion concerns tripling of the selected subset, not tripling of the whole original set.

Its strict-threshold partial structural statement sets $d=2|A|-|A^2|>0$ and defines ordered quotient multiplicities
\[
 r_{A^{-1}A}(g)=\#\{(a,b)\in A^2:a^{-1}b=g\},\qquad
 r_{AA^{-1}}(g)=\#\{(a,b)\in A^2:ab^{-1}=g\}.
\]
The aid claims
\[
 A^{-1}A=AA^{-1}=:Q,\qquad
 r_{A^{-1}A}(g)\geq d,\quad r_{AA^{-1}}(g)\geq d\quad(g\in Q),\qquad
 |Q|\leq |A|^2/d.
\]
It describes the strict threshold as sharp and points to later weak-Kneser and critical-pair theorems for broader inverse structure, without supplying their precise hypotheses or classification. Thus these explicit conclusions are recorded separately from the open-ended request for an inverse theorem.
\subsection{Short English statement}
Small doubling need not control tripling of the whole set in a nonabelian group. The supplied aid gives a large-subset replacement and a balanced-quotient statement below the strict doubling threshold two, while its broader inverse-theorem references remain unspecified.
\subsection{Sources}
\begin{itemize}
\item \emph{Problems Presented at the Workshop on Recent Trend in Additive Combinatorics}, collected by Ernie Croot and Vsevolod F. Lev, problem 3.5, T. Tao with the Ben Green/Mei-Chu Chang correction, page 8: \url{https://aimath.org/WWN/additivecomb/additivecomb.pdf}.
\item Local registry \texttt{data/unsolvedmath/problems.json}, permanent ID 20001117, exact source code \texttt{catalogue entry}; canonical location \texttt{aim-combinatorics-notes.json notes[241]}.
\item Attached attempt 1, partial result, problem status at run solved, rigor proved, low novelty confidence. The subset theorem, quotient lemma, and mentioned weak-Kneser/critical-pair results have no complete citations here. The aid is machine-generated, not peer reviewed, requiring independent verification before citation.
\end{itemize}
\end{document}
